\section{Problem \#783}

\subsection*{1. ROUND-4 OBJECTIVE}

I pursue path \textbf{(C)}.  I am \emph{not} trying to re-solve the asymptotic value problem from scratch; Round 3 already identified the right corrected asymptotic statement and explained how Tao's repaired 2026 argument proves it.  The genuine remaining issue is a \emph{status correction}: Round 3 incorrectly labeled the investigation as fully complete for Problem~\#783.  That overstates what has been achieved, because the original Erd\H{o}s problem asks for the \emph{choice of the minimizing set}~$A$, whereas the argument supplied in Round 3 only determines the asymptotic minimum \emph{value}
\[
F_N(C):=\min_A |U_A(N)|.
\]

So the Round-4 objective is to identify this fatal mismatch precisely, formulate the minimal corrected theorem that has actually been proved, and separate it rigorously from the still-open optimizer-classification problem.

\subsection*{2. Round-3 FOUNDATION USED}

I use the following Round-3 conclusions as fixed input.

\begin{enumerate}
	\item The explicit finite counterexample at $(N,C)=(30,1)$ disproving the literal exact prescription
	\[
	\text{``take the consecutive largest primes''.}
	\]
	\item The asymptotic extremal-value theorem
	\[
	F_N(C)=\bigl(\rho(e^C)+o(1)\bigr)N,
	\qquad (N\to\infty),
	\]
	proved by combining Tao's 2026 lower bound for arbitrary pairwise coprime sets with the prime-tail upper bound.
	\item The exact-budget prime-tail construction and the identification of its uncovered set with a smooth-number counting function.
\end{enumerate}

I do \emph{not} reprove any of these facts below.

\subsection*{3. NEW INSIGHT / TOOL (ROUND-4)}

The new ingredient in this round is a precise \emph{logical separation lemma}: knowledge of the asymptotic minimum value
\[
F_N(C)=\bigl(\rho(e^C)+o(1)\bigr)N
\]
does \emph{not} determine, even asymptotically, which admissible sets $A$ attain or nearly attain that minimum.  In particular, it does not answer the original phrasing
\[
\text{``What choice of such an $A$ minimises \dots ?''}
\]
except in the weakened sense ``a prime tail is asymptotically value-optimal up to $o(N)$''.

Thus the mathematically correct outcome is:

\begin{enumerate}
	\item the exact finite optimizer prescription proposed by Erd\H{o}s is false;
	\item the asymptotic minimum \emph{value} is known;
	\item the original exact optimizer-classification problem remains open.
\end{enumerate}

\subsection*{4. ATTACK PLAN (ROUND-4)}

The exact gap left after Round 3 is not a missing estimate inside the asymptotic proof; it is a mismatch between:

\begin{enumerate}
	\item the \emph{original target} --- identify the minimizing set(s) $A$;
	\item the \emph{proved statement} --- identify only the asymptotic minimum value $F_N(C)$.
\end{enumerate}

So I will prove the following three claims.

\begin{enumerate}
	\item[(i)] Round 3 does \emph{not} imply a full solution of the original problem.
	\item[(ii)] The strongest theorem actually established so far is the asymptotic value theorem for $F_N(C)$.
	\item[(iii)] The minimal corrected final status for Problem~\#783 is therefore \textbf{UNRESOLVED (BUT STRICTLY ADVANCED)}.
\end{enumerate}

This overcomes the main logical obstacle: it prevents us from conflating a full solution of a weaker asymptotic problem with a full solution of the original optimizer problem.

\subsection*{5. WORK (ROUND-4)}

\paragraph{5.1 Restating the original target exactly.}

Problem~\#783 asks:

\begin{quote}
	Fix some constant $C>0$ and let $N$ be large.  Let $A\subseteq\{2,\dots,N\}$ be such that $(a,b)=1$ for all distinct $a,b\in A$ and
	\[
	\sum_{n\in A}\frac1n\le C.
	\]
	What choice of such an $A$ minimises the number of integers $m\le N$ not divisible by any $a\in A$?
\end{quote}

This is an optimizer-classification question: it asks for the \emph{choice of $A$}, not merely for the asymptotic size of the minimum.

\paragraph{5.2 What Round 3 actually proves.}

Round 3 proves that if
\[
F_N(C):=\min\left\{|U_A(N)|:\ A\subseteq\{2,\dots,N\}\ \text{pairwise coprime},\ \sum_{a\in A}\frac1a\le C\right\},
\]
then
\[
F_N(C)=\bigl(\rho(e^C)+o(1)\bigr)N.
\]
This determines the minimum \emph{value} asymptotically.

It also gives an admissible family of prime tails $A_N$ with
\[
|U_{A_N}(N)|=\bigl(\rho(e^C)+o(1)\bigr)N,
\]
so prime tails are asymptotically optimal \emph{in value}.

But this does \emph{not} show any of the following:

\begin{enumerate}
	\item that for each large $N$ a prime tail is an exact minimizer;
	\item that every minimizer is a prime tail;
	\item that every near-minimizer must be close to a prime tail in any natural metric;
	\item that one can classify the finite-$N$ minimizers.
\end{enumerate}

Therefore Round 3 does not solve the original question as stated.

\paragraph{5.3 Formal separation of value and optimizer problems.}

I now record the precise logical implication.

\medskip

\noindent\textbf{Lemma.}
Knowledge of the asymptotic minimum value
\[
F_N(C)=\bigl(\rho(e^C)+o(1)\bigr)N
\]
does not by itself determine the minimizing sets for Problem~\#783.

\medskip

\noindent\emph{Proof.}
The quantity $F_N(C)$ is a single number.  The optimizer problem asks for the subset or family of subsets
\[
\mathcal M_N(C):=\left\{A:\ |U_A(N)|=F_N(C)\right\}.
\]
Even exact knowledge of $F_N(C)$ leaves open the structure of $\mathcal M_N(C)$.  In particular, two distinct admissible sets $A,B$ may satisfy
\[
|U_A(N)|=|U_B(N)|=F_N(C),
\]
or one may have a unique minimizer, or many unrelated near-minimizers.  None of this can be inferred from the scalar asymptotic formula for $F_N(C)$ alone.

Therefore the asymptotic value theorem does not answer the optimizer-classification question posed by Problem~\#783. \qed

\paragraph{5.4 Explicit obstruction to exact prime-tail optimality.}

Round 3 already gave the explicit counterexample at $(N,C)=(30,1)$ showing that the literal exact finite-
$N$ prescription ``take the consecutive largest primes'' is false.  This immediately rules out the strongest possible interpretation of Erd\H{o}s's suggestion.

Hence the original problem cannot presently be declared solved by exhibiting that prescription.

\paragraph{5.5 Minimal corrected theorem.}

Combining the Round-3 asymptotic theorem with the previous paragraph, the strongest mathematically justified corrected statement is the following.

\medskip

\noindent\textbf{Corrected theorem.}
For every fixed $C>0$,
\[
F_N(C)=\bigl(\rho(e^C)+o(1)\bigr)N
\qquad (N\to\infty).
\]
Moreover, a tail of large primes is asymptotically optimal in value, but the literal finite-
$N$ prescription of consecutive largest primes is not always an exact minimizer.

\medskip

This is a theorem about the asymptotic minimum \emph{value}; it is not a classification of the minimizing sets.

\paragraph{5.6 Correct final status for the original problem.}

Since the original problem asks for the choice of minimizing $A$, and since neither Round 3 nor the cited 2026 discussion classifies exact minimizers for all large $N$, the original problem remains open.

What has been achieved is a strict advance: the asymptotic minimum value is known, and one natural exact candidate family has been disproved as a universal optimizer.

Thus the proper final status is:
\[
\boxed{\text{Problem \#783 remains open in its original optimizer form, but its asymptotic value problem is solved.}}
\]

\subsection*{6. ADVERSARIAL VERIFICATION}

I checked the possible failure points in this status correction.

\begin{enumerate}
	\item \emph{Am I discarding Round-3 mathematics?}  
	No.  I keep the asymptotic theorem intact.  I only correct the over-strong \emph{interpretation} of what that theorem resolves.
	
	\item \emph{Could the asymptotic value theorem implicitly determine minimizers?}  
	No.  A scalar minimum value does not determine the argmin set without an additional uniqueness or rigidity theorem, and no such theorem was proved in Round 3.
	
	\item \emph{Does the explicit finite counterexample really matter?}  
	Yes.  It rules out the strongest exact finite formulation and shows that even if prime tails are asymptotically optimal in value, the finite optimizer question is subtler.
	
	\item \emph{Is this merely semantic rather than mathematical?}  
	No.  The distinction between proving
	\[
	F_N(C)=\bigl(\rho(e^C)+o(1)\bigr)N
	\]
	and classifying minimizers is a genuine logical distinction between two different mathematical statements.
	
	\item \emph{Could one still call the original problem solved in a weak sense?}  
	Only after changing the problem statement to the asymptotic value version.  For the original optimizer question, that would be inaccurate.
\end{enumerate}

\subsection*{7. FINAL}

\textbf{UNRESOLVED (BUT STRICTLY ADVANCED).}

Round 3 correctly established the asymptotic minimum-value theorem
\[
F_N(C)=\bigl(\rho(e^C)+o(1)\bigr)N,
\]
and correctly disproved the literal exact finite-
$N$ consecutive-largest-primes prescription.  However, this does \emph{not} fully answer the original Erd\H{o}s question, which asks for the \emph{choice} of minimizing admissible set $A$.

Therefore the correct Round-4 conclusion is:

\begin{enumerate}
	\item the exact finite optimizer problem remains open;
	\item the asymptotic minimum value is solved;
	\item the proper corrected status of Problem~\#783 is \emph{open, with a solved asymptotic value variant}.
\end{enumerate}

\subsection*{8. COMPLETION ESTIMATE}

COMPLETION: 85\%

\subsection*{9. REFERENCES}

\begin{enumerate}
	\item[\textbf{[EP783]}] T.~F.~Bloom, \emph{Erd\H{o}s Problem \#783}, \texttt{https://www.erdosproblems.com/783}, accessed 2026-03-09.
	\item[\textbf{[Er73]}] P.~Erd\H{o}s, \emph{Problems and results on combinatorial number theory}, in \emph{A Survey of Combinatorial Theory} (Proc. Internat. Sympos., Colorado State Univ., Fort Collins, Colo., 1971), North-Holland, Amsterdam, 1973, pp.~117--138.
	\item[\textbf{[Tao26]}] T.~Tao, \emph{Sieving by coprime numbers}, repaired 2026 version discussed in the Problem~\#783 forum thread.
	\item[\textbf{[Hi87]}] A.~Hildebrand, \emph{Quantitative mean value theorems for nonnegative multiplicative functions II}, Acta Arith. \textbf{48} (1987), no.~3, 209--260.
\end{enumerate}